% =====================================================================
%  Chapitre 6 — Conclusion et perspectives
%  Règles (Sayadi 5 / Buttler) :
%   - récapituler brièvement le cheminement et les conclusions intermédiaires
%   - énumérer les propositions
%   - terminer par limitations + perspectives futures
% =====================================================================
\chapter{Conclusion et perspectives}
\label{ch:conclusion}

% --- 6.1 Récapitulatif ----------------------------------------------
\section{Récapitulatif du travail réalisé}
\label{sec:concl:recap}

Ce mémoire a proposé une chaîne complète de détection de la fibrillation auriculaire à partir des seuls intervalles \RR{}, structurée en six phases.

La \textbf{phase~1} a établi un protocole d'évaluation patient-level rigoureux et caractérisé la forte hétérogénéité inter-patients de la base AFDB.

La \textbf{phase~2} a comparé quatre méthodes de référence (règle métier, forêt aléatoire, gradient boosting, \CNN{} 1D) sous ce même protocole.

La \textbf{phase~3} a conçu une architecture \CNN--\LSTM{} optimisée par recherche bayésienne, accompagnée d'une étude d'ablation.

La \textbf{phase~4} a évalué huit variantes de compression et identifié la combinaison optimale (quantification statique INT8) qui atteint 87~KB d'empreinte mémoire sans dégradation significative.

La \textbf{phase~5} a quantifié la perte cross-dataset par transfert AFDB~$\rightarrow$~LTAFDB.

La \textbf{phase~6} a livré un tableau de bord interactif \textit{Streamlit} pour l'exploration patient par patient des décisions du modèle.

% --- 6.2 Apports principaux -----------------------------------------
\section{Apports principaux}
\label{sec:concl:apports}

Les apports de ce travail sont les suivants :
\begin{itemize}[leftmargin=*]
  \item la mise en évidence d'un \keyword{plateau intrinsèque} de \(F_1 \approx 0{,}75\text{--}0{,}80\) pour la classification patient à partir des seuls \RR{}, indépendant de la famille de classifieur ;
  \item une étude de compression aboutie permettant un déploiement embarqué (modèle de 87~KB) ;
  \item une mesure honnête de la robustesse cross-dataset, peu documentée jusqu'ici ;
  \item un outil d'exploration interactive favorisant la transparence et la pédagogie auprès d'un public clinicien.
\end{itemize}

% --- 6.3 Limitations ------------------------------------------------
\section{Limitations}
\label{sec:concl:limites}

Les principales limitations du travail, déjà discutées au chapitre~\ref{ch:discussion}, sont rappelées ici :
\begin{itemize}
  \item un effectif patient restreint sur AFDB ;
  \item une dépendance aux annotations expertes pour les positions de R ;
  \item l'utilisation exclusive du signal \RR{}, qui ignore la morphologie ;
  \item un nombre limité de jeux de données externes pour l'évaluation de la robustesse.
\end{itemize}

% --- 6.4 Perspectives -----------------------------------------------
\section{Perspectives}
\label{sec:concl:perspectives}

Plusieurs extensions de ce travail sont envisagées.

\textbf{Court terme.}
\begin{itemize}
  \item Réaliser un déploiement de l'application sur Streamlit Community Cloud pour rendre la démonstration accessible publiquement.
  \item Compléter les ré-exports des figures en \texttt{.pdf} 300 dpi pour la qualité d'impression du présent mémoire.
\end{itemize}

\textbf{Moyen terme.}
\begin{itemize}
  \item Ajouter une seconde branche d'entrée portant la morphologie de l'\ECG{} (onde P, complexe QRS), afin de tester l'hypothèse du plateau intrinsèque.
  \item Évaluer le modèle final sur d'autres jeux de données publics (CPSC2018, Icentia11k) pour étendre la mesure de robustesse.
\end{itemize}

\textbf{Long terme.}
\begin{itemize}
  \item Industrialiser le pipeline avec un test prospectif sur un dispositif embarqué (microcontrôleur ARM Cortex-M, par exemple).
  \item Étudier une approche d'\emph{adaptation de domaine} (fine-tuning court par patient ou par dataset cible) pour combler l'écart cross-dataset.
\end{itemize}
